\chapter{Fragmentorbitalen en meeratomige moleculen}\label{ch:b2:fragment-orbitals}

De lewisstructuur van water toont twee gelijkwaardige O--H-bindingen en twee
gelijkwaardige vrije elektronenparen. Toch komen de elektronen, wanneer water door ultraviolet licht
geïoniseerd wordt, eruit met vier verschillende energieën, niet twee; methaan, met vier
gelijkwaardige C--H-bindingen, toont er twee. Elektronen in een molecuul zitten niet in bindingen
tussen paren atomen: ze bezetten orbitalen die over het hele molecuul uitgespreid zijn.
Dit hoofdstuk bouwt die orbitalen op uit de orbitalen van \omterm{def:b2:fragment-orbitals:fragment}{fragmenten}, verklaart
waarom water gebogen is en etheen vlak, en waarom een carbokation het liefst
alkylgroepen om zich heen heeft.

\begin{recall}
\Cref{ch:b2:diatomic-mos}: \omterm{def:b2:diatomic-mos:lcao}{LCAO}, overlap en symmetrie, twee wisselwerkende niveaus,
bindende en \omterm{def:b2:diatomic-mos:bonding}{antibindende orbitalen}. Het deel van jaar~1: VSEPR-vormen, hybridisatie
als beschrijving, resonantie, carbokationen en hun stabiliteitsvolgorde (de
wisselwerking van de C--H-bindingen met de lege orbitaal werd daar aangekondigd).
\end{recall}

\section{De fragmentmethode}

\begin{definition}[Fragmenten]\label{def:b2:fragment-orbitals:fragment}
Een \emph{fragment}\index{fragment} is een groep atomen van een molecuul die op
zichzelf beschouwd wordt (een atoom, een paar waterstofatomen, een \ce{CH2}-groep), en een
\emph{fragmentorbitaal}\index{fragmentorbitaal} is een van zijn orbitalen. De orbitalen van het
molecuul worden opgebouwd door de orbitalen van twee fragmenten twee aan twee met elkaar te laten
wisselwerken, zoals de atoomorbitalen van een tweeatomig molecuul.
\end{definition}

\begin{definition}[Symmetrie-aangepaste combinatie]\label{def:b2:fragment-orbitals:symmetry-adapted}
Een \emph{symmetrie-aangepaste combinatie}\index{symmetrie-aangepaste combinatie} van
gelijkwaardige atoomorbitalen is een combinatie die elke symmetrieoperatie van het
molecuul (een spiegeling in een spiegelvlak, een rotatie om een as) ofwel
onveranderd laat ofwel in haar tegengestelde omzet, of, voor ontaarde verzamelingen, omzet
in combinaties van dezelfde verzameling.
\end{definition}

\begin{proposition}[Twee waterstofatomen]\label{prop:b2:fragment-orbitals:h2-pair}
Voor twee gelijkwaardige waterstofatomen met 1s-orbitalen $h_1$ en $h_2$ zijn de
\omterm{def:b2:fragment-orbitals:symmetry-adapted}{symmetrie-aangepaste combinaties} $h_1 + h_2$ en $h_1 - h_2$: de eerste blijft
onveranderd, de tweede verandert van teken, onder elke operatie die de twee
atomen verwisselt.
\end{proposition}

\begin{proof}
Een operatie die de atomen verwisselt, zet $h_1$ om in $h_2$ en $h_2$ in
$h_1$: $h_1 + h_2$ blijft onveranderd en $h_1 - h_2$ wordt $h_2 - h_1 = -(h_1 - h_2)$.
\end{proof}

\begin{proposition}[Drie en vier waterstofatomen]\label{prop:b2:fragment-orbitals:h4}
Voor vier waterstofatomen op de hoekpunten van een tetraëder zijn de combinaties één
volledig symmetrische combinatie $h_1 + h_2 + h_3 + h_4$ en een verzameling van drie ontaarde
combinaties met elk één \omterm{def:b2:atomic-orbitals:node}{knoopvlak}, met de vorm van p-orbitalen. Voor drie
waterstofatomen op een driehoek (zoals in ammoniak): één symmetrische combinatie en een
ontaard paar.
\end{proposition}

\admitted

De vormen van de ontaarde combinaties, en hun namen (a, e, t), komen uit de
groepentheorie, die in het deel van jaar~3 ontwikkeld wordt; hier zijn de namen slechts labels.

\begin{proposition}[Alleen gelijk wisselwerkt met gelijk]\label{prop:b2:fragment-orbitals:interaction-rules}
Een \omterm{def:b2:fragment-orbitals:fragment}{fragmentorbitaal} wisselwerkt alleen met orbitalen van het andere \omterm{def:b2:fragment-orbitals:fragment}{fragment} die
zich onder elke symmetrieoperatie van het molecuul op dezelfde manier gedragen; de
sterkte van de wisselwerking groeit met de overlap en daalt met het
energieverschil.
\end{proposition}

\begin{proof}
Als twee orbitalen zich onder een of andere spiegeling verschillend gedragen, is de ene symmetrisch en
de andere antisymmetrisch ten opzichte van dat vlak, en zijn hun overlap en
\omterm{def:b2:diatomic-mos:integrals}{resonantie-integraal} nul (\cref{prop:b2:diatomic-mos:symmetry}). De afhankelijkheid van
overlap en energieverschil is die van \cref{thm:b2:diatomic-mos:heteronuclear}.
\end{proof}

\begin{method}[Een fragmentdiagram]\label{met:b2:fragment-orbitals:fragment-diagram}
\begin{enumerate}
  \item Splits het molecuul in twee \omterm{def:b2:fragment-orbitals:fragment}{fragmenten}, meestal een centraal atoom en de verzameling
        atomen eromheen.
  \item Bouw de \omterm{def:b2:fragment-orbitals:symmetry-adapted}{symmetrie-aangepaste combinaties} van de buitenste orbitalen op.
  \item Sorteer de orbitalen van beide \omterm{def:b2:fragment-orbitals:fragment}{fragmenten} naar hun gedrag onder de spiegelvlakken
        en assen van het molecuul.
  \item Laat elke orbitaal wisselwerken met die met hetzelfde gedrag, sterker
        naarmate ze dichter in energie liggen; wat overblijft, blijft niet-bindend.
  \item Vul met de valentie-elektronen en lees af: hoogste bezette niveau,
        bindend karakter, voorkeuren van vorm.
\end{enumerate}
\end{method}

\section{Water}

Plaats water in het $yz$-vlak, met de $z$-as als bissectrice van de hoek H--O--H. Twee
spiegelvlakken bevatten de $z$-as: het molecuulvlak $yz$ en het vlak $xz$
loodrecht erop. Ten opzichte van het vlak $xz$, dat de twee
waterstofatomen verwisselt, is $h_1 + h_2$ symmetrisch en $h_1 - h_2$ antisymmetrisch. Op zuurstof zijn 2s en
$2\mathrm p_z$ symmetrisch ten opzichte van beide vlakken; $2\mathrm p_y$ (in het
molecuulvlak, loodrecht op $z$) is antisymmetrisch ten opzichte van $xz$;
$2\mathrm p_x$ is antisymmetrisch ten opzichte van het molecuulvlak, wat geen enkele
waterstofcombinatie is.

Dus mengt $h_1 + h_2$ met 2s en $2\mathrm p_z$, $h_1 - h_2$ met $2\mathrm p_y$, en
blijft $2\mathrm p_x$ niet-bindend. Er ontstaan zes orbitalen, waarvan er vier gevuld worden door de acht
valentie-elektronen.

\begin{omfigure}
\begin{tikzpicture}[font=\small]
  % O fragment (left)
  \pic (os) at (0,-2.4) {omlevel={ud}{0.7}}; \node[left] at (-0.45,-2.4) {2s};
  \pic (op1) at (-0.45,0) {omlevel={ud}{0.4}}; \pic (op2) at (0,0) {omlevel={u}{0.4}}; \pic (op3) at (0.45,0) {omlevel={u}{0.4}};
  \node[left] at (-0.7,0) {2p};
  \node at (0,-3.2) {O};
  % H2 fragment (right)
  \pic (hp) at (6,0.9) {omlevel={u}{0.7}}; \node[right] at (6.4,0.9) {$h_1 + h_2$};
  \pic (hm) at (6,1.8) {omlevel={u}{0.7}}; \node[right] at (6.4,1.8) {$h_1 - h_2$};
  \node[anchor=north west] at (6.4,0.6) {fragment \ce{H2}};
  % MOs
  \pic (m1) at (3,-2.9) {omlevel={ud}{0.7}};  \node[omsymlabel, below=0.17cm, inner sep=0.5pt] at (m1-c) {$2a_1$};
  \pic (m2) at (3,-1.55) {omlevel={ud}{0.7}};  \node[omsymlabel, below=0.17cm, inner sep=0.5pt] at (m2-c) {$1b_2$};
  \pic (m3) at (3,-0.78) {omlevel={ud}{0.7}};  \node[omsymlabel, below=0.17cm, inner sep=0.5pt] at (m3-c) {$3a_1$};
  \pic (m4) at (3,0) {omlevel={ud}{0.7}};  \node[omsymlabel, below=0.17cm, inner sep=0.5pt] at (m4-c) {$1b_1$};
  \pic (m5) at (3,2.4) {omlevel={empty}{0.7}}; \node[omsymlabel, below=0.17cm, inner sep=0.5pt] at (m5-c) {$4a_1$};
  \pic (m6) at (3,3.1) {omlevel={empty}{0.7}}; \node[omsymlabel, below=0.17cm, inner sep=0.5pt] at (m6-c) {$2b_2$};
  \draw[omcorr, densely dashed] (os-r) -- (m1-l) (op3-r) -- (m2-l) (op3-r) -- (m3-l) (op3-r) -- (m4-l) (op3-r) -- (m5-l) (op3-r) -- (m6-l);
  \draw[omcorr, densely dashed] (hp-l) -- (m1-r) (hm-l) -- (m2-r) (hp-l) -- (m3-r) (hp-l) -- (m5-r) (hm-l) -- (m6-r);
\end{tikzpicture}

{\small Fragmentdiagram van water, schematisch. De symmetrische combinatie van de
waterstofatomen mengt met de 2s- en $2\mathrm p_z$-orbitalen van zuurstof ($a_1$-orbitalen),
de antisymmetrische met $2\mathrm p_y$ ($b_2$); $2\mathrm p_x$, loodrecht op het
molecuulvlak, blijft niet-bindend ($1b_1$, op het niveau van de 2p van zuurstof). Vier gevulde
valentieniveaus, vier foto-elektronbanden; de hoogste, $1b_1$, wordt geïoniseerd bij
\qty{12.62}{eV}. De labels
$a_1$, $b_1$, $b_2$ zijn namen (uit de groepentheorie, in het deel van jaar~3).}
% ledger: ie:H2O
\end{omfigure}

\begin{proposition}[Waarom water gebogen is]\label{prop:b2:fragment-orbitals:bent-water}
In een lineair H--O--H zijn twee p-orbitalen van zuurstof loodrecht op de as
niet-bindend. Buigen laat een ervan, die in het molecuulvlak, overlappen
met $h_1 + h_2$ en mengen met de 2s-orbitaal: dat gevulde niveau daalt, meer dan
de andere gevulde niveaus stijgen. Met acht valentie-elektronen is het gebogen molecuul
het stabielst.
\end{proposition}

\begin{proof}[Redenering]
In het lineaire molecuul ($z$ langs H--O--H) heeft $h_1 + h_2$ overlap nul met
$2\mathrm p_x$ en $2\mathrm p_y$ (elk is antisymmetrisch ten opzichte van een vlak dat
de as bevat, de combinatie symmetrisch); het zijn twee ontaarde
niet-bindende paren. Bij het buigen in het $yz$-vlak wijst $2\mathrm p_y$ gedeeltelijk naar
de waterstofatomen: haar overlap met $h_1 + h_2$ groeit vanaf nul, en het gevulde niveau
dat ze draagt, wordt gestabiliseerd (het niveau $3a_1$ van het gebogen molecuul). Het bindende
niveau $1b_2$ verliest een beetje overlap en stijgt licht. Het netto resultaat is een verlaging,
het grootst wanneer het dalende niveau gevuld is, zoals bij acht elektronen het geval is. Een
molecuul met slechts vier valentie-elektronen, zoals \ce{BeH2}, heeft dat niveau
leeg en is lineair.
\end{proof}

\section{Methaan, ammoniak en foto-elektronspectra}

In methaan geven de vier waterstoforbitalen de symmetrische combinatie, die
past bij de 2s-orbitaal van koolstof, en een ontaarde verzameling van drie, die past bij
zijn drie 2p-orbitalen. Vier \omterm{def:b2:diatomic-mos:bonding}{bindende orbitalen} zijn gevuld: één laag ($a_1$) en drie
ontaarde bij een hogere energie ($t_2$). Ammoniak, met drie waterstofatomen, houdt een
gevulde orbitaal op stikstof over die van de waterstofatomen weg wijst, zijn vrij elektronenpaar, als
hoogste bezette niveau.

\begin{omfigure}
\begin{tikzpicture}[font=\small]
  \pic (cs) at (0,-1.6) {omlevel={ud}{0.7}}; \node[left] at (-0.45,-1.6) {2s};
  \pic (cp1) at (-0.45,0) {omlevel={u}{0.4}}; \pic (cp2) at (0,0) {omlevel={u}{0.4}}; \pic (cp3) at (0.45,0) {omlevel={empty}{0.4}};
  \node[left] at (-0.7,0) {2p};
  \node at (0,-2.5) {C};
  \pic (ha) at (6,-0.8) {omlevel={ud}{0.7}}; \node[right] at (6.4,-0.8) {$a_1$};
  \pic (ht1) at (5.55,0.4) {omlevel={u}{0.4}}; \pic (ht2) at (6,0.4) {omlevel={u}{0.4}}; \pic (ht3) at (6.45,0.4) {omlevel={empty}{0.4}};
  \node[right] at (6.7,0.4) {$t_2$};
  \node at (6,-1.6) {fragment \ce{H4}};
  \pic (m1) at (3,-2.4) {omlevel={ud}{0.7}}; \node[omsymlabel, below=0.17cm, inner sep=0.5pt] at (m1-c) {$1a_1$};
  \pic (m2a) at (2.55,-0.9) {omlevel={ud}{0.4}}; \pic (m2b) at (3,-0.9) {omlevel={ud}{0.4}}; \pic (m2c) at (3.45,-0.9) {omlevel={ud}{0.4}};
  \node[omsymlabel, below=0.17cm, inner sep=0.5pt] at (m2b-c) {$1t_2$};
  \pic (m3) at (3,1.8) {omlevel={empty}{0.7}}; \node[omsymlabel, below=0.17cm, inner sep=0.5pt] at (m3-c) {$2a_1$};
  \pic (m4a) at (2.55,2.5) {omlevel={empty}{0.4}}; \pic (m4b) at (3,2.5) {omlevel={empty}{0.4}}; \pic (m4c) at (3.45,2.5) {omlevel={empty}{0.4}};
  \node[omsymlabel, below=0.17cm, inner sep=0.5pt] at (m4b-c) {$2t_2$};
  \draw[omcorr, densely dashed] (cs-r) -- (m1-l) (cp3-r) -- (m2a-l) (cs-r) -- (m3-l) (cp3-r) -- (m4a-l);
  \draw[omcorr, densely dashed] (ha-l) -- (m1-r) (ht1-l) -- (m2c-r) (ha-l) -- (m3-r) (ht1-l) -- (m4c-r);
\end{tikzpicture}

{\small Fragmentdiagram van methaan, schematisch. De symmetrische combinatie van de
vier waterstofatomen past bij de 2s-orbitaal van koolstof, de drie ontaarde bij zijn 2p-orbitalen:
twee gevulde niveaus, dus twee foto-elektronbanden, waarvan de bovenste
(het eerst geïoniseerd, bij \qty{12.61}{eV}) zes elektronen bevat.}
% ledger: ie:CH4
\end{omfigure}

\begin{definition}[Foto-elektronspectroscopie]\label{def:b2:fragment-orbitals:photoelectron}
\emph{Foto-elektronspectroscopie}\index{foto-elektronspectroscopie} meet de
kinetische energieën van de elektronen die uit moleculen worden weggeschoten door monochromatische
straling met bekende energie; elke band van het spectrum komt overeen met het verwijderen
van een elektron uit één bezet niveau, en de positie ervan geeft de ionisatie-energie
van dat niveau.
\end{definition}

\begin{proposition}[Benadering van Koopmans]\label{prop:b2:fragment-orbitals:koopmans}
De energie die nodig is om een elektron uit een bezette orbitaal te verwijderen, is
ongeveer min de energie van die orbitaal, $I \approx -\varepsilon$.
\end{proposition}

\admitted

De benadering verwaarloost de relaxatie van de andere elektronen na de
ionisatie; ze wordt in het deel van jaar~3 verantwoord. Ze maakt van een foto-elektronspectrum
een orbitaaldiagram: water toont vier valentiebanden, methaan twee,
in overeenstemming met de fragmentdiagrammen --- en in strijd met het beeld van twee gelijkwaardige
vrije elektronenparen en twee gelijkwaardige bindingen in water, dat dezelfde totale
elektronendichtheid op een andere manier beschrijft, maar niet de energieën geeft van de
verwijderde elektronen.

\section{Etheen}

Etheen kan opgebouwd worden uit twee \ce{CH2}-fragmenten. Elk \ce{CH2} heeft, naast de
orbitalen van zijn twee C--H-bindingen, een orbitaal die in het vlak van de waterstofatomen
weg wijst (die met haar partner de C--C-$\sigma$-binding vormt) en een p-orbitaal
loodrecht op het vlak. De twee p-orbitalen combineren tot een gevulde $\pi$ en een
lege $\pi^*$.

\begin{omfigure}
\begin{tikzpicture}[font=\small]
  \begin{scope}
    \pic at (0,0) {omp={0.85}{90}}; \pic at (1.2,0) {omp={0.85}{90}};
    \draw[thick] (0,0) -- (1.2,0);
    \draw[thick] (0,0) -- (-0.5,-0.45) (0,0) -- (-0.55,0.35) (1.2,0) -- (1.7,-0.45) (1.2,0) -- (1.75,0.35);
    \node at (0.6,-1.1) {$\pi$ (gevuld)};
  \end{scope}
  \begin{scope}[xshift=3.6cm]
    \pic at (0,0) {omp={0.85}{90}}; \pic at (1.2,0) {omp={-0.85}{90}};
    \draw[thick] (0,0) -- (1.2,0);
    \draw[thick] (0,0) -- (-0.5,-0.45) (0,0) -- (-0.55,0.35) (1.2,0) -- (1.7,-0.45) (1.2,0) -- (1.75,0.35);
    \node at (0.6,-1.1) {$\pi^*$ (leeg)};
  \end{scope}
  \begin{scope}[xshift=7.4cm]
    \pic at (0,0) {omp={0.85}{90}}; \pic at (1.2,0) {omp={0.85}{0}};
    \draw[thick] (0,0) -- (1.2,0);
    \node at (0.6,-1.1) {gedraaid over $90^\circ$};
    \node[font=\scriptsize, align=center] at (0.6,1.3) {overlap nul};
  \end{scope}
\end{tikzpicture}

{\small Het $\pi$-systeem van etheen, schematisch (het molecuul van opzij gezien, de
waterstofatomen in het vlak). De twee p-orbitalen loodrecht op het vlak combineren
in fase ($\pi$, gevuld) en in tegenfase ($\pi^*$, leeg). Eén \ce{CH2} over $90^\circ$
draaien maakt de twee p-orbitalen loodrecht op elkaar: hun overlap, en de $\pi$-binding,
verdwijnen.}
\end{omfigure}

De $\pi$-binding houdt etheen vlak: als één uiteinde gedraaid wordt, wijzen de p-orbitalen
van elkaar weg, daalt hun overlap als de cosinus van de draaihoek, en is bij $90^\circ$
de $\pi$-binding verdwenen. Om het molecuul te draaien moet de energie van een $\pi$-binding geleverd worden;
daarom zetten cis- en trans-isomeren van alkenen zich bij
kamertemperatuur niet in elkaar om, terwijl rotatie om een enkelvoudige binding vrij is.

\section{Hyperconjugatie}

\begin{definition}[Hyperconjugatie]\label{def:b2:fragment-orbitals:hyperconjugation}
\emph{Hyperconjugatie}\index{hyperconjugatie} is de wisselwerking van een gevulde
$\sigma$-orbitaal van een C--H- of C--C-binding met een naburige lege of gedeeltelijk gevulde
p- of $\pi^*$-orbitaal die er parallel aan ligt.
\end{definition}

\begin{proposition}[Alkylgroepen stabiliseren carbokationen]\label{prop:b2:fragment-orbitals:hyperconjugation}
Een carbokation wordt gestabiliseerd door elke C--H- of C--C-binding op een naburig koolstofatoom die
zich kan richten naar zijn lege p-orbitaal: hoe meer zulke bindingen, hoe stabieler het
kation, vandaar de volgorde tertiair > secundair > primair > methyl.
\end{proposition}

\begin{proof}[Redenering]
De lege p-orbitaal van het kationische koolstofatoom en een gevulde orbitaal $\sigma_{\text{C--H}}$
ernaast vormen een systeem met twee niveaus, met een groot energieverschil $\Delta$ en een kleine koppeling
$\beta$ (alleen verschillend van nul als de binding ongeveer parallel ligt aan de p-orbitaal). Volgens
\cref{thm:b2:diatomic-mos:heteronuclear} wordt het gevulde $\sigma$-niveau verlaagd met ongeveer
$\beta^2/\Delta$, en winnen de twee elektronen het dubbele daarvan; het lege niveau stijgt, zonder
kosten. Elke goed gerichte binding voegt haar eigen stabilisatie toe; een methylgroep
heeft altijd één C--H-binding die bijna goed gericht is, wat haar rotatie ook is.
\end{proof}

\begin{omfigure}
\begin{tikzpicture}[font=\small]
  \pic at (0,0) {omp={0.9}{90}};
  \node[circle, fill=black, inner sep=1.2pt] at (0,0) {};
  \node[font=\scriptsize, anchor=north west] at (0.08,-0.05) {C$^+$};
  \draw[thick] (0,0) -- (1.4,0);
  \node[circle, fill=black, inner sep=1.2pt] at (1.4,0) {};
  \draw[thick] (1.4,0) -- (1.75,0.85);
  \pic at (1.58,0.42) {omlobe={0.75}{68}};
  \pic at (1.58,0.42) {omlobe={0.45}{248}};
  \node[font=\scriptsize] at (1.95,1.0) {H};
  \draw[thick] (1.4,0) -- (2.0,-0.45) (1.4,0) -- (1.1,-0.6);
  \draw[<->, omProp, thick] (0.35,0.8) .. controls (0.8,1.2) and (1.2,1.2) .. (1.5,0.95);
  \node[font=\scriptsize, align=left, anchor=west] at (2.4,0.4) {gevulde lob $\sigma_{\text{C--H}}$ parallel\\ aan de lege p-orbitaal};
  % level scheme
  \begin{scope}[xshift=7.2cm, yshift=-0.4cm]
    \pic (p) at (0,1.4) {omlevel={empty}{0.7}}; \node[left] at (-0.4,1.4) {lege p};
    \pic (s) at (2.2,-0.8) {omlevel={ud}{0.7}}; \node[right] at (2.6,-0.8) {$\sigma_{\text{C--H}}$};
    \pic (lo) at (1.1,-1.2) {omlevel={ud}{0.7}};
    \pic (hi) at (1.1,1.7) {omlevel={empty}{0.7}};
    \draw[omcorr, densely dashed] (p-r) -- (lo-l) (p-r) -- (hi-l) (s-l) -- (lo-r) (s-l) -- (hi-r);
    \node[font=\scriptsize, anchor=north] at (1.1,-1.35) {gestabiliseerd met $\approx \beta^2/\Delta$};
  \end{scope}
\end{tikzpicture}

{\small \omterm{def:b2:fragment-orbitals:hyperconjugation}{Hyperconjugatie} in een carbokation, schematisch. Een gevulde bindende C--H-orbitaal
op het naburige koolstofatoom, parallel aan de lege p-orbitaal, staat een deel van zijn
elektronendichtheid eraan af; het gevulde niveau wordt gestabiliseerd zoals bij elke wisselwerking van
twee niveaus met verschillende energie.}
\end{omfigure}

\section{Oefeningen}

\begin{exercise}[$\star$]\label{exo:b2:fragment-orbitals:1}
Schrijf de twee \omterm{def:b2:fragment-orbitals:symmetry-adapted}{symmetrie-aangepaste combinaties} op van de 1s-orbitalen van de waterstofatomen van
water en zeg welke ervan van teken verandert in het vlak dat de hoek H--O--H middendoor deelt,
loodrecht op het molecuul.
\end{exercise}

\begin{exercise}[$\star$]\label{exo:b2:fragment-orbitals:2}
Hoeveel \omterm{def:b2:diatomic-mos:lcao}{valentie-molecuulorbitalen} heeft water, hoeveel valentie-elektronen,
en hoeveel gevulde valentieorbitalen?
\end{exercise}

\begin{exercise}[$\star$]\label{exo:b2:fragment-orbitals:3}
Waarom toont het foto-elektronspectrum van methaan twee valentiebanden, waarvan de ene
driemaal zo intens is als de andere?
\end{exercise}

\begin{exercise}[$\star$]\label{exo:b2:fragment-orbitals:4}
Waarom liggen de zes atomen van etheen in één vlak?
\end{exercise}

\begin{exercise}[$\star\star$]\label{exo:b2:fragment-orbitals:5}
De eerste ionisatie-energie van ammoniak is \qty{10.07}{eV}, die van methaan
\qty{12.61}{eV}. Welke orbitaal verliest in elk geval het elektron, en waarom is ammoniak
gemakkelijker te ioniseren?
\end{exercise}

\begin{exercise}[$\star\star$]\label{exo:b2:fragment-orbitals:6}
Vergelijk de bezette niveaus van lineair en gebogen water, en leg uit waarom een molecuul
zoals \ce{BeH2}, met vier valentie-elektronen, lineair is.
\end{exercise}

\begin{exercise}[$\star\star$]\label{exo:b2:fragment-orbitals:7}
Verklaar met het resultaat voor twee niveaus de stabiliteitsvolgorde van de carbokationen
\ce{CH3+}, \ce{CH3CH2+}, \ce{(CH3)2CH+}, \ce{(CH3)3C+}.
\end{exercise}

\begin{exercise}[$\star\star$]\label{exo:b2:fragment-orbitals:8}
De overlap van twee p-orbitalen die over een hoek $\theta$ om de binding gedraaid zijn, is
evenredig met $\cos\theta$. Hoe hangt de $\pi$-stabilisatie af van $\theta$, als de
\omterm{def:b2:diatomic-mos:integrals}{resonantie-integraal} evenredig is met de overlap? Bij welke hoek is ze tot de helft
gedaald?
\end{exercise}

\begin{exercise}[$\star\star$]\label{exo:b2:fragment-orbitals:9}
Boraan \ce{BH3} is vlak, ammoniak piramidaal. Verklaar het verschil met de \omterm{def:b2:fragment-orbitals:fragment}{fragmentorbitalen} van een
centraal atoom en drie waterstofatomen, aan de hand van het aantal elektronen van elk.
\end{exercise}

\begin{exercise}[$\star\star\star$]\label{exo:b2:fragment-orbitals:10}
Het ion \ce{H3+} is een gelijkzijdige driehoek. Bepaal met $\alpha$ en $\beta$ voor de 1s-orbitalen,
en met verwaarloosde overlap, zijn drie orbitaalenergieën (de symmetrische
combinatie heeft energie $\alpha + 2\beta$) en leg uit waarom twee elektronen volstaan om
het te binden.
\end{exercise}

\begin{exercise}[$\star\star\star$]\label{exo:b2:fragment-orbitals:11}
Het allylsysteem \ce{CH2=CH-CH2} heeft drie parallelle p-orbitalen. Raad uit het aantal knopen de vormen
van zijn drie $\pi$-orbitalen (tekens op elk koolstofatoom), en welke ervan
het hoogste bezette niveau van het allylanion is.
\end{exercise}

\begin{exercise}[$\star\star\star$]\label{exo:b2:fragment-orbitals:12}
Bouw het $\pi$-systeem van koolstofdioxide op uit de 2p-orbitalen loodrecht op de
as O--C--O: hoeveel $\pi$-orbitalen, welke zijn bindend, niet-bindend,
antibindend, en hoeveel elektronen bevatten ze?
\end{exercise}

\section{Probleem: de vorm van water}

\begin{problem}[{Weekendprobleem --- de \omterm{def:b2:fragment-orbitals:fragment}{fragmenten} zuurstof en \ce{H2}, de orbitalen van
een lineair H--O--H, de niveaus die verschuiven wanneer het buigt, en het foto-elektronspectrum
met de benadering van Koopmans}]
\label{pb:b2:fragment-orbitals:1}
Gegevens: geometrie van water, $r(\ce{O-H}) = \qty{95.8}{pm}$, hoek
$104.48^\circ$; eerste ionisatie-energieën: \ce{H2O} \qty{12.62}{eV}, \ce{O}
\qty{13.62}{eV}. Water ligt in het $yz$-vlak, met $z$ langs de bissectrice.
% ledger: geo:H2O.rOH, geo:H2O.aHOH, ie:H2O, ie:O

\medskip
\noindent\textbf{Deel I --- \omterm{def:b2:fragment-orbitals:fragment}{Fragmenten}.}
\begin{enumerate}
  \item Hoeveel valentie-elektronen heeft water, en van welke atomen komen ze?
  \item Schrijf de twee \omterm{def:b2:fragment-orbitals:symmetry-adapted}{symmetrie-aangepaste combinaties} van de 1s-orbitalen van de waterstofatomen op.
  \item Welke orbitalen van zuurstof zijn symmetrisch ten opzichte van beide spiegelvlakken?
  \item Welke orbitaal van zuurstof is alleen antisymmetrisch ten opzichte van het vlak $xz$?
  \item Welke orbitaal van zuurstof heeft geen partner onder de waterstofcombinaties?
  \item Hoeveel \omterm{def:b2:diatomic-mos:lcao}{molecuulorbitalen} ontstaan er, en hoeveel daarvan zijn gevuld?
\end{enumerate}

\medskip
\noindent\textbf{Deel II --- Een lineair H--O--H.}
\begin{enumerate}[resume]
  \item Welke combinatie wisselwerkt in een lineair molecuul langs $z$ met de
        $2\mathrm p_z$ van zuurstof?
  \item Welke wisselwerkt met de 2s van zuurstof?
  \item Welke orbitalen van zuurstof blijven niet-bindend, en met welke ontaarding?
  \item Vul de niveaus met de acht valentie-elektronen.
  \item Zou lineair water \omterm{def:b2:diatomic-mos:magnetism}{paramagnetisch} zijn?
  \item Waar zou zijn hoogste bezette niveau liggen, vergeleken met een 2p-orbitaal
        van zuurstof?
\end{enumerate}

\medskip
\noindent\textbf{Deel III --- Buigen.}
\begin{enumerate}[resume]
  \item Het lineaire molecuul langs $z$ wordt gebogen in het $yz$-vlak. Welke
        \omterm{def:b2:diatomic-mos:bonding}{niet-bindende orbitaal} begint te overlappen met $h_1 + h_2$?
  \item Wat gebeurt er met het gevulde niveau dat ze draagt?
  \item Welk bindend niveau stijgt licht, en waarom?
  \item Besluit: waarom is water gebogen?
  \item Vergelijk de gemeten hoek met $90^\circ$ (zuivere p-binding) en $109.5^\circ$.
  \item Welke \omterm{def:b2:diatomic-mos:bonding}{niet-bindende orbitaal} blijft door het buigen onaangeroerd?
\end{enumerate}

\medskip
\noindent\textbf{Deel IV --- Het foto-elektronspectrum.}
\begin{enumerate}[resume]
  \item Hoeveel valentiebanden toont het spectrum?
  \item Uit welke orbitaal komt volgens de benadering van Koopmans de eerste band?
  \item Waarom is die band smal, terwijl banden van \omterm{def:b2:diatomic-mos:bonding}{bindende orbitalen} breed zijn?
  \item Vergelijk de eerste ionisatie-energie van water met die van het zuurstofatoom,
        en verklaar het verschil.
  \item Wat wordt er van de ``twee gelijkwaardige vrije elektronenparen'' van de lewisstructuur?
  \item Geef het aantal valentie-ionisatiebanden van water en de energie van
        de laagste.
\end{enumerate}
\end{problem}
